% =============================================================================
% Monthly control variables in the GCHP CO2 4D-Var system:
% what the adjoint accumulator means and how to extract dJ/dsigma_m
% =============================================================================
\documentclass[12pt]{article}

% Latin Modern ships Type 1 outlines; the plain Computer Modern path needs
% MetaFont (mf), which is absent on this host and aborts the build.
\usepackage{lmodern}
\usepackage{amsmath,amssymb,amsfonts}
\usepackage{geometry}
\usepackage{hyperref}
\usepackage{booktabs}
\usepackage{xcolor}
\newcommand{\bm}[1]{\mathbf{#1}}

\geometry{margin=1in}

\newcommand{\Fprior}{F^{\mathrm{p}}}
\newcommand{\tend}{t_{\mathrm{end}}}
\newcommand{\tstart}{t_{\mathrm{start}}}
\newcommand{\Adj}{A}

\title{\textbf{Monthly Control Variables in the GCHP CO$_2$ 4D-Var System}\\
\large Why the monthly gradient is a difference of the adjoint accumulator}
\author{}
\date{\today}

\begin{document}
\maketitle

\begin{abstract}
Extending the control vector from a single annual scaling field $\sigma$ to twelve
monthly fields $\sigma_1,\dots,\sigma_{12}$ requires knowing what the adjoint
diagnostic \texttt{SurfaceFluxAdj\_CO2} actually contains. This note shows that the
diagnostic is a \emph{backward accumulator in emission time}, that the monthly
gradient is therefore the difference of the accumulator between consecutive month
boundaries, and that using the raw accumulator value at each month boundary
multiplies the true gradient by a lower-triangular matrix of ones. It also
separates two ideas that are easy to conflate: influence in \emph{observation
time} (a December observation does inform January fluxes, and this is retained)
versus contribution in \emph{emission time} (which is what the differencing
partitions). A source-code audit confirms that the accumulator is weighted by the
\emph{prior} TER flux, so the differenced gradient needs no $1/\sigma_m$
renormalisation. A finite-difference verification protocol is given.
\end{abstract}

% -----------------------------------------------------------------------------
\section{Setup and notation}
% -----------------------------------------------------------------------------

The assimilation window is $[\tstart,\tend]$, here 2016-01-01 to 2017-01-01.
Let $t_1 < t_2 < \dots < t_{12}$ be the month start times, $t_1 = \tstart$, and
define the closing boundary $t_{13} = \tend$. Month $m$ is the interval
$[t_m, t_{m+1})$, and $m(s)$ denotes the month containing time $s$.

The control variable multiplies the terrestrial ecosystem respiration (TER) prior
flux. Writing $\Fprior(x,s)$ for the prior TER flux at surface location $x$ and
time $s$, the flux actually seen by the model is
%
\begin{equation}
  F_\sigma(x,s) \;=\; \sigma_{m(s)}(x)\,\Fprior(x,s).
  \label{eq:flux}
\end{equation}
%
This is implemented as HEMCO scale factor \texttt{750} (\texttt{Sigma\_CO2}) on the
\texttt{TER\_CLASS\_CTEM} base entry. Note that $\sigma$ scales TER \emph{only}:
fossil fuel, ocean, CARDAMOM NBE and GPP are not scaled by it.

The cost function is $J = J_{\mathrm{obs}} + J_{\mathrm{b}}$, where $J_{\mathrm{obs}}$
sums observation-space misfits over all observation times $t' \in [\tstart,\tend]$
and $J_{\mathrm{b}}$ is the background term.

Let $\lambda(x,s)$ denote the sensitivity of $J_{\mathrm{obs}}$ to an
\emph{instantaneous} surface flux input at $(x,s)$,
%
\begin{equation}
  \lambda(x,s) \;\equiv\; \frac{\partial J_{\mathrm{obs}}}{\partial F(x,s)} .
  \label{eq:lambda}
\end{equation}
%
This is the quantity the adjoint transport carries; it is the surface trace of the
adjoint tracer field \texttt{SpeciesAdj\_CO2}.

% -----------------------------------------------------------------------------
\section{What the adjoint carries backward}
% -----------------------------------------------------------------------------

The adjoint model is integrated backward from $\tend$ with $\lambda \equiv 0$ as its
terminal condition, and observation forcing is injected as the integration passes
each observation time. Consequently, at any time $s$,
%
\begin{equation}
  \lambda(x,s) \;=\; \sum_{\substack{\text{observations at } t' \\ t' \,\ge\, s}}
    \Big( \text{adjoint transport of that observation's forcing from } t' \text{ back to } s \Big).
  \label{eq:lambda_sum}
\end{equation}
%
\textbf{This is the key structural fact.} For $s$ in January, $\lambda(x,s)$ already
contains the forcing of \emph{every} observation from January through December. A
January flux perturbation needs months to reach the mid-troposphere where XCO$_2$
is sensitive; the adjoint propagates that December information backward in time
and deposits it into $\lambda(x,s)$ for $s$ in January. Nothing in what follows
discards it.

% -----------------------------------------------------------------------------
\section{The diagnostic is an accumulator in emission time}
% -----------------------------------------------------------------------------

The GCHP adjoint does not archive $\lambda$ at the surface directly. It archives a
running sum, taken over the backward sweep, of $\lambda$ weighted by the prior flux.
Writing $\Adj(x,t)$ for \texttt{SurfaceFluxAdj\_CO2},
%
\begin{equation}
  \boxed{\;\Adj(x,t) \;=\; \sum_{s \,\ge\, t} \lambda(x,s)\,\Fprior(x,s)\,\Delta t \;}
  \label{eq:accumulator}
\end{equation}
%
with the sum running over model timesteps $s$ from $t$ up to $\tend$.

Two immediate consequences, both of which are observed in the archived output of the
existing one-month OSSE run:
%
\begin{enumerate}
  \item $\Adj(x,\tend) = 0$ identically (empty sum). Verified: the final record of
        \texttt{GEOSChem.Adjoint.20160201.nc4} is exactly zero.
  \item $\Adj(x,\tstart)$ is the complete window-integrated sensitivity. This is the
        value that \texttt{lbfgsb\_step.py} currently reads, and it is the correct
        gradient \emph{for a single scaling field applied over the whole window}.
\end{enumerate}

Indeed, if $\sigma$ is one field applied to all of $[\tstart,\tend]$, then by
\eqref{eq:flux} and the chain rule
%
\begin{equation}
  \frac{\partial J_{\mathrm{obs}}}{\partial \sigma(x)}
  \;=\; \sum_{s \in [\tstart,\tend]} \lambda(x,s)\,\Fprior(x,s)\,\Delta t
  \;=\; \Adj(x,\tstart),
\end{equation}
%
which is exactly what the current one-month script does.

% -----------------------------------------------------------------------------
\section{The monthly gradient}
% -----------------------------------------------------------------------------

Now let $\sigma_m$ act on month $m$ alone. Since $\sigma_m$ multiplies $\Fprior(x,s)$
if and only if $s \in [t_m, t_{m+1})$, the chain rule gives
%
\begin{equation}
  g_m(x) \;\equiv\; \frac{\partial J_{\mathrm{obs}}}{\partial \sigma_m(x)}
  \;=\; \sum_{s \in [t_m,\,t_{m+1})} \lambda(x,s)\,\Fprior(x,s)\,\Delta t .
  \label{eq:gm}
\end{equation}
%
Comparing \eqref{eq:gm} with the accumulator \eqref{eq:accumulator}, the sum over
month $m$ is the difference of two partial sums:
%
\begin{equation}
  \boxed{\;g_m(x) \;=\; \Adj(x,t_m) \;-\; \Adj(x,t_{m+1}),
  \qquad \Adj(x,t_{13}) = \Adj(x,\tend) = 0. \;}
  \label{eq:diff}
\end{equation}

\paragraph{Telescoping check.} Summing \eqref{eq:diff} over $m=1,\dots,12$ telescopes:
$\sum_m g_m = \Adj(\tstart) - \Adj(\tend) = \Adj(\tstart)$. The twelve monthly
gradients must add up to the single annual gradient. This is a cheap and strong
runtime assertion. On the existing one-month output, splitting the month into four
sub-windows and summing the four differences reproduces $\Adj(\tstart)$ to
$3.8 \times 10^{-6}$, i.e.\ to \texttt{float32} roundoff on a field of magnitude
$\sim\!2.9\times10^{3}$.

% -----------------------------------------------------------------------------
\section{Why the raw accumulator value is not the monthly gradient}
% -----------------------------------------------------------------------------

The tempting alternative is to read $\Adj(x,t_m)$ directly and call it $g_m$. From
\eqref{eq:accumulator} and \eqref{eq:gm},
%
\begin{equation}
  \Adj(x,t_m) \;=\; \sum_{k=m}^{12} g_k(x),
  \label{eq:nested}
\end{equation}
%
so the raw values are \emph{nested partial sums} of the true gradients. In matrix
form, with $\bm{g} = (g_1,\dots,g_{12})^{\!\top}$ and
$\bm{g}^{\mathrm{raw}} = (\Adj(t_1),\dots,\Adj(t_{12}))^{\!\top}$,
%
\begin{equation}
  \bm{g}^{\mathrm{raw}} \;=\; \bm{U}\,\bm{g},
  \qquad
  U_{mk} = \begin{cases} 1 & k \ge m \\ 0 & k < m \end{cases}
\end{equation}
%
an upper-triangular matrix of ones. In particular $\Adj(t_1) = \sum_k g_k$: the raw
January value is the \emph{whole year's} gradient, and $\Adj(t_{12}) = g_{12}$ is the
only raw value that happens to be correct.

The error is not subtle. It says that scaling January's TER changes $J$ by an amount
that includes the effect of scaling February's, March's, \dots, December's TER --- even
though $\sigma_1$ has no influence whatsoever on an emission that occurs in December.

\paragraph{Directional-derivative test.} For a perturbation
$\bm{\delta} = (\delta_1,\dots,\delta_{12})$ the true first-order change is
$\delta J = \sum_k g_k \delta_k$. Using the raw values instead gives
%
\begin{equation}
  \sum_m \Adj(t_m)\,\delta_m
  \;=\; \sum_m \delta_m \sum_{k \ge m} g_k
  \;=\; \sum_k g_k \Big( \sum_{m \le k} \delta_m \Big),
\end{equation}
%
which equals $\delta J$ only when $\delta_m = 0$ for all $m > 1$. For a uniform
perturbation $\delta_m = \delta$ it returns $\delta \sum_k k\, g_k$: month $k$ is
weighted by $k$, so December is counted twelve times and January once. L-BFGS-B fed
this quantity would systematically overstate the early months.

% -----------------------------------------------------------------------------
\section{Observation time versus emission time}
% -----------------------------------------------------------------------------

It is worth stating plainly what the differencing does and does not remove, because
the two sums in play are indexed by different variables.

\begin{itemize}
  \item \textbf{Observation time $t'$.} This is summed inside $\lambda$, equation
        \eqref{eq:lambda_sum}. For $s$ in January the sum runs over all $t' \ge s$,
        i.e.\ over January through December observations. \emph{Differencing never
        touches this sum.} The physical statement ``a December XCO$_2$ observation
        informs January surface fluxes'' is encoded here, and it survives intact.

  \item \textbf{Emission time $s$.} This is summed in the accumulator, equation
        \eqref{eq:accumulator}. Differencing restricts this sum to
        $s \in [t_m, t_{m+1})$, because $\sigma_m$ multiplies nothing outside its own
        month.
\end{itemize}

So $g_{\mathrm{Jan}} = \Adj(t_{\mathrm{Jan}}) - \Adj(t_{\mathrm{Feb}})$ retains the
influence of December \emph{observations} while excluding the contribution of
February--December \emph{emissions}. Both are required, and they are not in tension.

\paragraph{A real, separate limitation.} There \emph{is} a genuine finite-window
problem nearby, and it should not be confused with the above. Fluxes emitted in
November and December have little time to influence any observation before the
window closes at $\tend$, so by \eqref{eq:gm} their gradients $g_{11}, g_{12}$ are
small; the background term dominates and those months remain close to their prior.
This is a property of a one-year window and appears under \emph{any} correct
gradient convention. Differencing does not cause it --- it merely makes it legible.
Mitigations, if the late months matter scientifically, are to extend $\tend$ a few
months past the last month being optimised, or to accept that the last two monthly
increments are prior-dominated.

% -----------------------------------------------------------------------------
\section{Normalisation: the accumulator uses the prior flux}
% -----------------------------------------------------------------------------

Equation \eqref{eq:accumulator} asserts that the accumulator weights $\lambda$ by the
\emph{prior} flux $\Fprior$, not by the current flux $F_\sigma = \sigma\,\Fprior$. Had
the code weighted by $F_\sigma$, we would instead have
%
\begin{equation}
  \tilde{\Adj}(t_m) - \tilde{\Adj}(t_{m+1}) \;=\; \sigma_m(x)\, g_m(x),
  \label{eq:norm}
\end{equation}
%
requiring an explicit division by $\sigma_m$. The two coincide at $\sigma \equiv 1$, so
the distinction is invisible at iteration~1 and would corrupt every later iteration.

\paragraph{This is settled by the source code, and the prior flux is used.} The
accumulation is performed in \texttt{Adjoint\_Utils\_Mod.F90}:
%
\begin{verbatim}
State_Chm%SurfaceFluxAdj(:,:,NFD) = State_Chm%SurfaceFluxAdj(:,:,NFD) +
      State_Chm%SpeciesAdj(:,:,1,NFD) *
      ( SurfaceFluxForward * dt / ( rho_dry * dz ) )
\end{verbatim}
%
where \texttt{SurfaceFluxForward} is \texttt{State\_Chm\%SurfaceFlux}, populated in the
adjoint run by \texttt{Compute\_Sflx\_For\_Adjoint} (\texttt{hco\_interface\_gc\_mod.F90})
according to the selector in the adjoint \texttt{GCHP.rc}:
%
\begin{verbatim}
ADJ_HEMCO_SFLUX_SELECTOR: DIAGN
ADJ_HEMCO_SFLUX_DIAGN:    EmisCO2_NetTerrExch
\end{verbatim}
%
Crucially, in the \emph{adjoint} \texttt{HEMCO\_Config.rc} the \texttt{TER\_CLASS\_CTEM}
base entry carries scale factor \texttt{751} (the unit conversion) \emph{only}, whereas
the \emph{forward} \texttt{HEMCO\_Config.rc} carries \texttt{750/751}, where \texttt{750}
is \texttt{Sigma\_CO2}. The adjoint therefore evaluates \texttt{EmisCO2\_NetTerrExch}
as the \emph{unscaled prior} TER flux at every timestep, for every iteration. This is
the asymmetry that makes \eqref{eq:accumulator} --- and hence \eqref{eq:diff} --- correct
as written, with no $1/\sigma_m$ factor.

\paragraph{Corollary.} The adjoint never reads \texttt{Sigma\_CO2} in any way that
affects its output. Consequently \textbf{only the forward control files require the
monthly sigma plumbing}; the adjoint \texttt{ExtData.rc} / \texttt{HEMCO\_Config.rc}
need no change for the monthly extension.

\paragraph{What this does \emph{not} explain.} The unphysical growth of $\sigma$
(max $\approx 6.2$) observed in the one-month OSSE at iterations $\ge 2$ is therefore
\emph{not} attributable to a $\sigma$-weighted accumulator. It remains open, and the
finite-difference protocol of Section~\ref{sec:verify} is still worth running, both to
confirm the differencing convention and to bound the gradient error independently of
this argument.

% -----------------------------------------------------------------------------
\section{Verification protocol}
\label{sec:verify}
% -----------------------------------------------------------------------------

The following finite-difference test discriminates the two remaining candidate
conventions (differencing vs.\ raw) at the cost of a few short forward runs, and
independently re-checks the $\Fprior$ normalisation established in Section~6. Use a
two-month window, $\tstart = $ 2016-01-01, $\tend = $ 2016-03-01, so that
$m \in \{1,2\}$.

\begin{enumerate}
  \item Run forward + adjoint once with $\sigma_1 = \sigma_2 = 1$. Record $J^{(0)}$ and
        the accumulator at $t_1, t_2, t_3$. Form
        $g_2^{\mathrm{diff}} = \Adj(t_2) - \Adj(t_3) = \Adj(t_2)$ and
        $g_2^{\mathrm{raw}} = \Adj(t_2)$. (These coincide for the \emph{last} month ---
        see \eqref{eq:nested} --- so month~2 tests only the normalisation.)
  \item Form $g_1^{\mathrm{diff}} = \Adj(t_1) - \Adj(t_2)$ and
        $g_1^{\mathrm{raw}} = \Adj(t_1)$. These differ, so month~1 discriminates the
        conventions.
  \item Choose a single grid cell $x_\star$ with a large $|g_1|$. Rerun the
        \emph{forward model only} with $\sigma_1(x_\star) = 1 + \varepsilon$,
        $\varepsilon = 10^{-2}$, all else unity. Record $J^{(1)}$.
  \item Compare
        $\displaystyle \frac{J^{(1)} - J^{(0)}}{\varepsilon}$
        against $g_1^{\mathrm{diff}}(x_\star)$ and $g_1^{\mathrm{raw}}(x_\star)$.
        The differencing convention should match to $O(\varepsilon)$; the raw value
        should not.
  \item \textbf{Normalisation re-check.} Repeat step~3 about a non-unit base point: set
        $\sigma_1(x_\star) = 2$ for the reference run and $\sigma_1(x_\star) = 2+\varepsilon$
        for the perturbed run, rerunning the adjoint at the base point. Per Section~6 the
        recovered finite difference should equal the $\Adj$-difference \emph{directly}
        (prior-flux weighting). If it instead equals the $\Adj$-difference divided by $2$,
        the code is accumulating $F_\sigma$ (equation \eqref{eq:norm}), the source audit
        is wrong somewhere, and \texttt{lbfgsb\_step\_monthly.py} must divide by $\sigma_m$.
\end{enumerate}

A cheaper partial check, requiring no extra runs, is the telescoping assertion
$\sum_{m} \big[\Adj(t_m) - \Adj(t_{m+1})\big] = \Adj(\tstart)$, which should be added
to \texttt{lbfgsb\_step\_monthly.py} as a runtime guard.

% -----------------------------------------------------------------------------
\section{Summary}
% -----------------------------------------------------------------------------

\begin{center}
\begin{tabular}{@{}ll@{}}
\toprule
Quantity & Expression \\
\midrule
Archived diagnostic & $\Adj(x,t) = \sum_{s \ge t} \lambda(x,s)\Fprior(x,s)\Delta t$ \\
Weight (from source) & prior TER flux, \texttt{EmisCO2\_NetTerrExch}, no \texttt{750} \\
Terminal condition & $\Adj(x,\tend) = 0$ \\
Annual gradient (current code) & $\partial J_{\mathrm{obs}}/\partial\sigma = \Adj(x,\tstart)$ \\
Monthly gradient (new code) & $g_m(x) = \Adj(x,t_m) - \Adj(x,t_{m+1})$ \\
Runtime guard & $\sum_{m=1}^{12} g_m = \Adj(x,\tstart)$ \\
Full gradient & $\nabla_{\sigma_m} J = g_m + (\sigma_m - 1)/\sigma_b^2$ \\
\bottomrule
\end{tabular}
\end{center}

\noindent
The December-observation argument is correct physics and is fully retained by
\eqref{eq:diff}; it constrains $\lambda$, not the range of the emission-time sum.
The accumulated weight is the prior flux $\Fprior$, established by source-code audit in
Section~6, so \eqref{eq:diff} is used as written and no $\sigma_m$ renormalisation is
needed. What remains is to confirm the convention numerically via Section~\ref{sec:verify}.

\end{document}
